\textbf{Phần 1. Trắc nghiệm lựa chọn}
\setcounterpageref{ex}{0}

\begin{ex} %[3D2V26]
Thực hiện phép tính: $84 : 4 + 19$. Kết quả là
\choice
{$38$}
{$42$}
{$36$}
{\True $40$}
\loigiai{
$84 : 4 = 21$; $21 + 19 = 40$.
\textbf{Đáp án của câu hỏi là D}
}
\end{ex}

\begin{ex} %[3D2V26]
Một cuộn dây thép dài $96\text{ m}$ được cắt thành 3 đoạn bằng nhau. Độ dài mỗi đoạn dây là
\choice
{$28\text{ m}$}
{\True $32\text{ m}$}
{$36\text{ m}$}
{$30\text{ m}$}
\loigiai{
Độ dài mỗi đoạn dây là: $96 : 3 = 32\text{ (m)}$.
\textbf{Đáp án của câu hỏi là B}
}
\end{ex}

\begin{ex} %[3D2V26]
Tìm $x$, biết: $5 \times x = 95 - 20$.
\choice
{$13$}
{$17$}
{\True $15$}
{$19$}
\loigiai{
$5 \times x = 75 \Rightarrow x = 75 : 5 = 15$.
\textbf{Đáp án của câu hỏi là C}
}
\end{ex}

\begin{ex} %[3D2V26]
Có 72 học sinh xếp đều thành 6 hàng. Hỏi mỗi hàng có bao nhiêu học sinh?
\choice
{\True $12$ học sinh}
{$14$ học sinh}
{$10$ học sinh}
{$16$ học sinh}
\loigiai{
Mỗi hàng có số học sinh là: $72 : 6 = 12$ (học sinh).
\textbf{Đáp án của câu hỏi là A}
}
\end{ex}

\begin{ex} %[3D2V26]
Bác thợ may dùng một tấm vải dài $68\text{ m}$ để may các bộ rèm cửa, mỗi bộ rèm cửa cần $4\text{ m}$ vải. Bác đã may được 12 bộ rèm. Hỏi bác còn có thể may thêm được bao nhiêu bộ rèm cửa nữa?
\choice
{$4$ bộ rèm}
{$6$ bộ rèm}
{$7$ bộ rèm}
{\True $5$ bộ rèm}
\loigiai{
Số bộ rèm may được từ tấm vải là: $68 : 4 = 17$ (bộ). Bác còn may được thêm: $17 - 12 = 5$ (bộ rèm).
\textbf{Đáp án của câu hỏi là D}
}
\end{ex}

\begin{ex} %[3D2V26]
Thương của phép chia số lớn nhất có hai chữ số cho số chẵn lớn nhất có một chữ số là
\choice
{$11$ (dư $7$)}
{\True $12$ (dư $3$)}
{$12$ (dư $5$)}
{$10$ (dư $9$)}
\loigiai{
Số lớn nhất có 2 chữ số là 99. Số chẵn lớn nhất có 1 chữ số là 8. Ta có: $99 : 8 = 12$ (dư 3).
\textbf{Đáp án của câu hỏi là B}
}
\end{ex}

\begin{ex} %[3D2V26]
Một cửa hàng có 90 bao gạo. Sau khi bán đi một số bao gạo thì số bao gạo còn lại bằng $\dfrac{1}{5}$ số bao gạo ban đầu. Hỏi cửa hàng đã bán đi bao nhiêu bao gạo?
\choice
{$18$ bao gạo}
{$80$ bao gạo}
{\True $72$ bao gạo}
{$64$ bao gạo}
\loigiai{
Số bao gạo còn lại: $90 : 5 = 18$ (bao). Số bao gạo đã bán: $90 - 18 = 72$ (bao gạo).
\textbf{Đáp án của câu hỏi là C}
}
\end{ex}

\begin{ex} %[3D2V26]
Tìm số tự nhiên $y$ biết: $y \times 4 + 18 = 74$.
\choice
{\True $14$}
{$16$}
{$12$}
{$15$}
\loigiai{
$y \times 4 = 74 - 18 = 56 \Rightarrow y = 56 : 4 = 14$.
\textbf{Đáp án của câu hỏi là A}
}
\end{ex}

\begin{ex} %[3D2V26]
Người ta đóng 78 lon nước ngọt vào các lốc, mỗi lốc có 6 lon. Hỏi đóng được tất cả bao nhiêu lốc nước ngọt?
\choice
{$12$ lốc}
{$14$ lốc}
{$15$ lốc}
{\True $13$ lốc}
\loigiai{
Số lốc nước ngọt đóng được là: $78 : 6 = 13$ (lốc).
\textbf{Đáp án của câu hỏi là D}
}
\end{ex}

\begin{ex} %[3D2V26]
Một sợi dây dài $85\text{ dm}$ được chia làm 5 đoạn bằng nhau. Độ dài 3 đoạn dây như thế là
\choice
{$17\text{ dm}$}
{\True $51\text{ dm}$}
{$34\text{ dm}$}
{$68\text{ dm}$}
\loigiai{
Một đoạn dài: $85 : 5 = 17\text{ (dm)}$. Ba đoạn dài: $17 \times 3 = 51\text{ (dm)}$.
\textbf{Đáp án của câu hỏi là B}
}
\end{ex}

\begin{ex} %[3D2V26]
Bác Năm thu hoạch được $92\text{ kg}$ cà chua, bác chia đều vào 4 sọt. Mỗi sọt bác lại đem bán bớt đi $8\text{ kg}$. Hỏi mỗi sọt còn lại bao nhiêu ki-lô-gam cà chua?
\choice
{$23\text{ kg}$}
{$18\text{ kg}$}
{\True $15\text{ kg}$}
{$12\text{ kg}$}
\loigiai{
Mỗi sọt ban đầu có: $92 : 4 = 23\text{ (kg)}$. Còn lại: $23 - 8 = 15\text{ (kg)}$.
\textbf{Đáp án của câu hỏi là C}
}
\end{ex}

\begin{ex} %[3D2V26]
Một hình thoi có chu vi bằng $68\text{ cm}$. Biết hình thoi có 4 cạnh dài bằng nhau. Độ dài mỗi cạnh của hình thoi là
\choice
{\True $17\text{ cm}$}
{$16\text{ cm}$}
{$18\text{ cm}$}
{$15\text{ cm}$}
\loigiai{
Độ dài mỗi cạnh là: $68 : 4 = 17\text{ (cm)}$.
\textbf{Đáp án của câu hỏi là A}
}
\end{ex}

\textbf{Phần 2. Trắc nghiệm đúng sai}
\setcounterpageref{ex}{0}

\begin{ex} %[3D2V26]
Xét tính đúng sai của các phép chia số có hai chữ số cho số có một chữ số:
\choiceTF
{\True Khi chia số có hai chữ số cho số có một chữ số, ta thực hiện chia lần lượt từ trái sang phải.}
{\True Kết quả của $78 : 3$ là $26$.}
{Phép chia $85 : 5$ có kết quả là $15$.}
{\True Trong phép chia $91 : 7$, thương nhận được là một số có hai chữ số.}
\loigiai{
\begin{itemize}
    \item Ý a) Đúng: Quy tắc chia theo thứ tự hàng chục rồi hàng đơn vị.
    \item Ý b) Đúng: $78 : 3 = 26$.
    \item Ý c) Sai: Sai vì $85 : 5 = 17$.
    \item Ý d) Đúng: $91 : 7 = 13$, là số có 2 chữ số.
\end{itemize}
\begin{itemize}
    \item a) \textbf{Đúng}.
    \item b) \textbf{Đúng}.
    \item c) \textbf{Sai}.
    \item d) \textbf{Đúng}.
\end{itemize}
}
\end{ex}

\begin{ex} %[3D2V26]
Một xưởng bánh có 96 chiếc bánh mì nướng được đóng vào các hộp 4 chiếc hoặc 6 chiếc:
\choiceTF
{\True Nếu đóng vào các hộp 4 chiếc thì đóng được 24 hộp.}
{\True Nếu đóng vào các hộp 6 chiếc thì đóng được 16 hộp.}
{Số hộp 4 chiếc ít hơn số hộp 6 chiếc là 8 hộp.}
{\True Cả hai cách đóng đều chia hết, không thừa chiếc bánh nào.}
\loigiai{
\begin{itemize}
    \item Ý a) Đúng: $96 : 4 = 24$ hộp.
    \item Ý b) Đúng: $96 : 6 = 16$ hộp.
    \item Ý c) Sai: Sai vì số hộp 4 chiếc nhiều hơn: $24 - 16 = 8$ hộp.
    \item Ý d) Đúng: 96 chia hết cho cả 4 và 6.
\end{itemize}
\begin{itemize}
    \item a) \textbf{Đúng}.
    \item b) \textbf{Đúng}.
    \item c) \textbf{Sai}.
    \item d) \textbf{Đúng}.
\end{itemize}
}
\end{ex}

\textbf{Phần 3. Trả lời ngắn}
\setcounterpageref{ex}{0}

\begin{ex} %[3D2V26]
Tính kết quả của phép chia: $96 : 4$.
\shortans{24}
\loigiai{
\textbf{Đáp số:} $24$\\
$96 : 4 = 24$.
}
\end{ex}

\begin{ex} %[3D2V26]
Một cửa hàng có 65 bao xi măng chia đều cho 5 công trình. Hỏi mỗi công trình nhận được bao nhiêu bao xi măng?
\shortans{13}
\loigiai{
\textbf{Đáp số:} $13$\\
$65 : 5 = 13$ bao xi măng.
}
\end{ex}

\begin{ex} %[3D2V26]
Tìm giá trị của $x$, biết: $x \times 6 = 78$.
\shortans{13}
\loigiai{
\textbf{Đáp số:} $13$\\
$x = 78 : 6 = 13$.
}
\end{ex}

\begin{ex} %[3D2V26]
Một lớp học có 48 học sinh, trong đó $\dfrac{1}{3}$ số học sinh là học sinh giỏi. Hỏi lớp đó có bao nhiêu học sinh giỏi?
\shortans{16}
\loigiai{
\textbf{Đáp số:} $16$\\
Số học sinh giỏi là: $48 : 3 = 16$ học sinh.
}
\end{ex}

\begin{ex} %[3D2V26]
Đoạn đường từ nhà Nam đến trường dài $84\text{ m}$. Nam đã đi được $\dfrac{1}{4}$ quãng đường. Hỏi Nam còn phải đi bao nhiêu mét nữa để đến trường?
\shortans{63}
\loigiai{
\textbf{Đáp số:} $63$\\
Nam đã đi: $84 : 4 = 21\text{ m}$. Nam còn phải đi: $84 - 21 = 63\text{ (m)}$.
}
\end{ex}

\begin{ex} %[3D2V26]
Tìm số tự nhiên có một chữ số $k$ sao cho: $98 : k = 14$.
\shortans{7}
\loigiai{
\textbf{Đáp số:} $7$\\
$k = 98 : 14 = 7$.
}
\end{ex}
